Digital messaging platforms such as WhatsApp offer limited support for modifying or withdrawing access to shared photos. Existing controls offer only partial protection and do not support ongoing, revisable consent at the level of an individual photo. This project investigates how consent revocation can be designed as a usable interaction within WhatsApp-style photo sharing. Drawing on literature concerning dynamic consent, multi-party privacy, and privacy-focused HCI, four research gaps are identified. In response, \textbf{Smart Photos} is proposed which is a permission model where each shared photo carries a persistent and revisable permission state covering forwarding, downloading, view-once, time limited access, and co-subject governance. An interactive front-end prototype was built with React, Next.js and TypeScript and evaluated through semi-structured interviews with 9 participants using scenario-based prototype materials. Reflexive thematic analysis used and five themes are extracted which shows that participants valued revisable permissions as a meaningful improvement while recognising that revocation remains inherently partial. The dissertation also discusses recommendations and future directions for different stakeholder groups that seek to adopt a similar model. 